\chapter{Entropic OT, Sinkhorn, and Barycentric Projection}
\label{ch:bf-dpf-entropic-ot-sinkhorn}

The previous chapter fixed the deterministic OT equal-weighting baseline.  This
chapter explains the main differentiable computational route used in the OT lane:
entropic regularization, finite Sinkhorn iteration, barycentric projection, and
the distinction between the mathematical teacher object and the numerical object
actually differentiated in code.

This chapter has one dominant goal: a reader should be able to understand what
object finite Sinkhorn computes, how barycentric output is produced from it, where
approximation enters, and what diagnostics must be checked before the result is
used in a downstream filtering scalar.

\section{What changes when entropy is added}
\label{sec:bf-eot-what-changes}

The deterministic OT baseline chooses a coupling that minimizes transport cost
exactly for the declared linear program.  Entropic OT changes that object by
adding a smoothing term to the optimization problem.  In plain language, the
transport plan is now rewarded for being more diffuse instead of concentrating all
mass as sharply as possible.

The gain is numerical smoothness and a tractable matrix-scaling solver.  The cost
is that the transport object is no longer the unregularized deterministic OT
solution.  This is the first approximation layer of the chapter.

\section{Running three-particle cell with entropic smoothing}
\label{sec:bf-eot-running-cell}

Keep the same running cell as before:
\[
  x_1=-2,
  \qquad x_2=0,
  \qquad x_3=3,
  \qquad
  w=
  \left(
    \frac{1}{2},
    \frac{1}{3},
    \frac{1}{6}
  \right),
  \qquad
  u=
  \left(
    \frac{1}{3},
    \frac{1}{3},
    \frac{1}{3}
  \right),
\]
with quadratic cost matrix
\begin{equation}
\label{eq:bf-eot-running-cost}
  C
  =
  \begin{pmatrix}
    0 & 4 & 25 \\
    4 & 0 & 9 \\
    25 & 9 & 0
  \end{pmatrix}.
\end{equation}
For \(\varepsilon=1\), the Gibbs kernel is
\begin{equation}
\label{eq:bf-eot-running-kernel}
  K_1=
  \begin{pmatrix}
    1 & e^{-4} & e^{-25} \\
    e^{-4} & 1 & e^{-9} \\
    e^{-25} & e^{-9} & 1
  \end{pmatrix}.
\end{equation}
The entries are heavily concentrated near the diagonal, but not exactly zero away
from it.  That is a simple visual sign of what entropy has changed: the plan is no
longer forced into the sharpest possible assignment.

\section{Entropic OT teacher object}
\label{sec:bf-eot-teacher-object}

Let \(\mathcal U(w,u)\) be the feasible coupling set from the previous chapter.
The entropically regularized OT problem is
\begin{equation}
\label{eq:bf-eot-primal}
  \Pi_\varepsilon^\star
  =
  \arg\min_{\Pi\in\mathcal U(w,u)}
  \left\{
    \langle \Pi,C\rangle
    + \varepsilon\sum_{i,j}\Pi_{ij}(\log\Pi_{ij}-1)
  \right\}.
\end{equation}
This chapter uses the same entropy convention as the earlier OT material, so the
minimizer is the exact teacher for the chosen regularized object, not for the
unregularized linear program and not for categorical resampling.

\paragraph{Reader checkpoint.}
At this stage, the chapter has changed the mathematical object once: from the
unregularized deterministic OT coupling to the entropically smoothed OT coupling.
Nothing numerical has happened yet.  The next step explains how this exact
regularized object is computed in practice.

\subsection{Exact teacher versus exact engineering route}
\label{sec:bf-eot-teacher-versus-engineering}

The scalable-OT survey suggests a useful extra distinction before moving on to
finite Sinkhorn iterations.  There are two different senses in which a route may
still be called ``exact'' for the entropic OT object.

First, there is the mathematical teacher object
\(\Pi_\varepsilon^\star\) from \eqref{eq:bf-eot-primal}.  That is the exact
minimizer of the entropically regularized problem.  Second, there is an
engineering route that preserves that same entropic object while changing how the
matrix-scaling updates are executed in memory or on hardware.  Exact online,
streaming, and GPU Sinkhorn methods belong to this second category: they do not
change the entropic teacher problem, but they avoid materializing the full cost,
kernel, or coupling whenever possible.

This distinction matters for the BayesFilter transport lane.  A method that
recomputes the same entropic scaling updates through blocked reductions or
streamed kernel application is still on the semantics-preserving reference lane.
By contrast, a method that replaces the kernel, truncates the coupling family, or
changes the transport objective has moved to a different approximation class even
if it is motivated by scalability.  The reader should therefore keep separate:
\begin{enumerate}[label=(\roman*)]
  \item the exact entropic teacher object,
  \item the finite numerical object produced by a declared Sinkhorn execution
    path,
  \item the engineering strategy used to realize that numerical path.
\end{enumerate}
The next sections fix the second object explicitly and later staged chapters will
discuss scalability routes that either preserve or alter it.

\section{Why the optimizer factorizes}
\label{sec:bf-eot-factorization}

The first-order optimality conditions imply that the regularized coupling has the
form
\begin{equation}
\label{eq:bf-eot-factorization}
  \Pi_\varepsilon^\star
  =
  \diag(a)K_\varepsilon\diag(b),
  \qquad
  (K_\varepsilon)_{ij}=\exp(-C_{ij}/\varepsilon),
\end{equation}
with positive scaling vectors \(a,b\) chosen so that the coupling has the correct
row and column sums.  Operationally, this means the hard OT problem has been
reduced to the problem of finding two scaling vectors that balance a known kernel.
That reduction is why Sinkhorn iteration is so useful.

\section{Algorithm: finite Sinkhorn equal-weighting}
\label{sec:bf-eot-algorithm}

\begin{algorithm}[ht]
\caption{\texttt{sinkhorn\_equalweight\_step}}
\begin{algorithmic}[1]
\Require weighted particle cloud \(\{x_i,w_i\}_{i=1}^N\), equal target marginal \(u\), cost matrix \(C\), regularization \(\varepsilon>0\), initialization \((a^{(0)},b^{(0)})\), iteration budget \(K\), and stabilization policy
\Ensure finite coupling \(\Pi_{\varepsilon,K}\), barycentric equal-weight cloud \(\widetilde X\), row/column residuals, and declared numerical metadata
\State Form the Gibbs kernel \((K_\varepsilon)_{ij}=\exp(-C_{ij}/\varepsilon)\)
\For{\(k=0,\ldots,K-1\)}
  \State Update \(a^{(k+1)} = w \oslash (K_\varepsilon b^{(k)})\)
  \State Update \(b^{(k+1)} = u \oslash (K_\varepsilon' a^{(k+1)})\)
\EndFor
\State Form the finite coupling \(\Pi_{\varepsilon,K}=\diag(a^{(K)})K_\varepsilon\diag(b^{(K)})\)
\State Compute residuals
\[
  r_{\mathrm{row}}=\|\Pi_{\varepsilon,K}\mathbf 1-w\|_1,
  \qquad
  r_{\mathrm{col}}=\|\Pi_{\varepsilon,K}'\mathbf 1-u\|_1
\]
\For{each output slot \(j=1,\ldots,N\)}
  \State Compute \(\tilde x^{(j)} = \frac{1}{u_j}\sum_{i=1}^N \Pi_{ij}x_i\)
\EndFor
\State Return the coupling, equal-weight cloud, residuals, \(\varepsilon\), \(K\), initialization policy, and stabilization metadata
\end{algorithmic}
\end{algorithm}

\subsection{Exact online and streaming Sinkhorn as the reference lane}
\label{sec:bf-eot-reference-lane}

The scalable-OT survey also clarifies how BayesFilter should talk about scaling
without silently changing the transport object.  A finite Sinkhorn route can be
made substantially more practical by changing the execution strategy rather than
the underlying entropic teacher problem.  Streaming, online, or GPU-specialized
implementations still belong to the reference lane when they preserve the same
entropic scaling equations and the same barycentric output rule while avoiding
full materialization of the cost or kernel matrix.

In that reference-lane setting, the implementation may compute the needed
row-wise or column-wise reductions in blocks, keep only partial transport state
in memory, or fuse barycentric application with solver passes.  What must remain
fixed is the declared numerical object: the finite coupling
\(\Pi_{\varepsilon,K}\), its residual checks, and the barycentric output derived
from that same finite coupling.  Thus a streamed or GPU-tiled Sinkhorn routine is
not automatically a new approximation class.  It becomes a new approximation
class only when the scaling equations, kernel, feasible family, or output object
itself are changed.

This chapter therefore treats exact-online and streaming implementations as
engineering variants of the same entropic reference lane.  Later staged chapters
on retained-teacher approximations will take up the different case where the
transport object itself is compressed, factorized, or otherwise altered for
scalability.

This is the chapter’s main callable procedure.  A reader who implements only one
OT route from the monograph should be able to implement this one.

\section{Worked balancing pass on the running cell}
\label{sec:bf-eot-worked-balance}

Take the simplest initialization \(b^{(0)}=(1,1,1)'\).  Then the first
row-balancing step computes
\begin{equation}
\label{eq:bf-eot-first-row-balance}
  a^{(1)} = w \oslash (K_1 b^{(0)}).
\end{equation}
For the running cell,
\[
  K_1 b^{(0)}
  \approx
  (1.0183,\,1.0184,\,1.0001),
\]
so
\[
  a^{(1)}
  \approx
  (0.491,\,0.327,\,0.167).
\]
The next column-balancing step then uses these row-repaired values to push the
column sums toward \((1/3,1/3,1/3)\).  The exact arithmetic is less important
than the interpretation: finite Sinkhorn repeatedly repairs whichever marginal is
currently wrong.

\paragraph{Plain-language takeaway.}
Sinkhorn is not a mysterious OT oracle.  It is an alternating balancing procedure
for a smoothed transport matrix.  That is why later learned warm-start methods try
to predict better initial balancing coordinates rather than replacing the solver
outright.

\section{Barycentric projection as the output map}
\label{sec:bf-eot-barycentric}

The solver output is still only a coupling matrix.  The equal-weight particle cloud
that later code consumes is obtained by barycentric projection:
\begin{equation}
\label{eq:bf-eot-barycentric}
  \tilde x^{(j)}
  =
  \frac{1}{u_j}\sum_{i=1}^N \Pi_{ij}x_i.
\end{equation}
In the equal-weight case, this is
\begin{equation}
\label{eq:bf-eot-barycentric-equal}
  \tilde x^{(j)}=N\sum_{i=1}^N \Pi_{ij}x_i.
\end{equation}
This step is where the numerical coupling becomes the actual equal-weight output
cloud.  It is therefore the place where an implementation must be especially clear
about shapes, normalization, and any auxiliary state that is or is not carried
alongside the particle locations.

\section{Differentiation contracts}
\label{sec:bf-eot-differentiation}

The final major distinction of the chapter is between three different derivative
objects:
\begin{enumerate}[label=(\roman*)]
  \item derivative of the exact regularized OT optimizer,
  \item derivative of the finite unrolled Sinkhorn computation,
  \item derivative of some scalar built from that finite computation.
\end{enumerate}
In code, the second and third are the most common.  A gradient is only the
correct gradient of the exact executed scalar if the computational graph used in
autodiff matches the scalar value actually reported by the code.

\section{Filtering-loop insertion point}
\label{sec:bf-eot-loop}

Inside a filter, the finite EOT/Sinkhorn equal-weighting step sits after the
weighted particle cloud has been formed and before an equally weighted cloud is
passed onward:
\begin{equation}
\label{eq:bf-eot-loop}
  \{x_i,w_i\}_{i=1}^N
  \xrightarrow{\mathrm{EOT/Sinkhorn}}
  \{\tilde x^{(j)},1/N\}_{j=1}^N.
\end{equation}
If the broader filtering algorithm carries auxiliary particle-local state, this
chapter does not yet assert how that state is transported.  Its contract is
narrower: it specifies the equal-weighting rule for the particle locations and the
mass-routing matrix from which that equal-weight cloud is derived.

\section{Verification contract}
\label{sec:bf-eot-verification}

A reader implementing this route should verify at least:
\begin{enumerate}[label=(\arabic*)]
  \item row and column residuals are reported and decrease with additional
  iterations or stronger stabilization,
  \item the initialization, \(\varepsilon\), iteration budget, and stabilization
  policy are recorded,
  \item the barycentric cloud has the declared shape and normalization,
  \item any reported gradient is tied to the exact computational graph that was
  actually executed,
  \item any stronger filtering or HMC claim is kept separate from these local
  numerical checks.
\end{enumerate}

A failed residual check is evidence about the finite solver or stabilization path,
not about the abstract EOT object.  A successful residual check is evidence that
the numerical route is behaving as declared, not that it preserves the original
categorical target.

\section{What this chapter does and does not claim}
\label{sec:bf-eot-boundary}

This chapter teaches the main differentiable OT computational route in a near-
implementation style.  It does \emph{not} claim:
\begin{itemize}
  \item that entropic OT is the same object as unregularized OT,
  \item that finite Sinkhorn is the exact regularized optimizer,
  \item that a barycentric equal-weight cloud is the same object as categorical
  ancestor sampling,
  \item that autodiff through the executed graph is automatically the gradient of
  the original filtering likelihood.
\end{itemize}
Its role is to fix the entropic teacher object, the finite numerical route, the
output cloud construction, and the differentiation contract that the recovered
custom-gradient bridge and later learned-transport chapters build on.

\FloatBarrier
